\section{Μαθηματική Ανακατασκευή Σημάτων}
\subsection{Εφαρμογή της Μεθόδου Matrix Pencil (MP)}
Η ανάλυση των αποκρίσεων \textit{ringdown} μέσω τεχνικών ταυτοποίησης συστήματος αποτελεί μια μέθοδο για την παρακολούθηση της δυναμικής ευστάθειας \cite{modal2018}. Η διαδικασία της ταυτοποίησης βασίστηκε στην ανάλυση των προεπεξεργασμένων σημάτων $V$, $I$, $P$ και $Q$ μέσω του αλγορίθμου MP. Σύμφωνα με τις απαιτήσεις της διερεύνησης, η μέθοδος εφαρμόστηκε με δύο διαφορετικές προσεγγίσεις για τον καθορισμό της τάξης $n$ του συστήματος. 

\begin{enumerate}
    \item \textbf{Καθορισμένη Τάξη:} Χρησιμοποιήθηκαν σταθερές τιμές τάξης $n \in \{2, 4, 6\}$, ενώ για το \textit{IEEE 39-Bus System} εξετάστηκε επιπλέον και η τιμή $n=8$, καθώς πρόκειται για σύστημα με πλουσιότερο modal περιεχόμενο και μεγαλύτερο πλήθος ηλεκτρομηχανικών ρυθμών σε σύγκριση με το \textit{Kundur Two-Area System}~\cite{sarkar2024drga}
    \item \textbf{Αυτόματος Καθορισμός Τάξης:} Εφαρμόστηκε η τεχνική της αποκοπής των ιδιοτιμών με παραμέτρους ανοχής $\tau \in \{1, 0.1, 0.01\}$, ε\-πι\-τρέ\-πο\-ντας στον αλγόριθμο να επιλέξει αυτόματα την τάξη του μοντέλου με βάση τη διαδοχική βελτίωση της ανακατασκευής του σήματος. Στο πλαίσιο αυτό, η τάξη αυξανόταν μέχρι η μεταβολή του συντελεστή  $R^2$  μεταξύ δύο διαδοχικών ανακατασκευών να γίνει μικρότερη απο $\tau$.
\end{enumerate}

\subsection{Μαθηματική Ανακατασκευή και Πιστότητα}
Για κάθε σενάριο ταυτοποίησης εξήχθησαν οι παράμετροι των ιδιοτιμών: το πλάτος ($A_{i}$), ο ρυθμός απόσβεσης ($\sigma_{i}$), η συχνότητα ($f_{i}$ σε Hz) και η φάση ($\varphi_{i}$). Η μαθηματική ανακατασκευή του σήματος $\hat{y}(t)$ πραγματοποιήθηκε μέσω της σχέσης:

\begin{equation}
\hat{y}(t)=\sum_{i=1}^{n}A_{i}e^{\sigma_{i}t}\cos(2\pi f_{i}t+\varphi_{i})
\label{eq:reconstruction}
\end{equation}

Η πιστότητα της ανακατασκευής αξιολογήθηκε μέσω του συντελεστή προσδιορισμού $R^{2}$ και του μέσου τετραγωνικού σφάλματος RMSE. Η διαδικασία αυτή εξασφαλίζει ότι οι ταυτοποιημένες ιδιοτιμές αντιπροσωπεύουν με ακρίβεια τη φυσική απόκριση των εξεταζόμενων συστημάτων και δεν αποτελούν μαθηματικά σφάλματα της προσέγγισης.


